% Part: computability
% Chapter: computability-theory
% Section: halting-problem

\documentclass[../../../include/open-logic-section]{subfiles}

\begin{document}

\olfileid{cmp}{thy}{hlt}
\olsection{حساب رُکّن دا مسئلہ}

تعمیر مطابق، ہمہ گیر جزوی قابلِ حساب فنکشن~$\fn{Un}(e,x)$
اُدوں تے صرف اُدوں تعریف شدہ اے جدوں~$e$ دے کوڈ والا فنکشن
انپٹ~$x$ اُتے حساب کرکے کوئی قدر دیوے۔ فطری سوال ایہ اے کہ
کیہ اسی فیصلہ کر سکدے آں کہ ایہ حالت اے یا نہیں۔ اصل وچ
ایہ فیصلہ نہیں کیتا جا سکدا۔ حساب دے ٹیورنگ مشین نمونے وچ
ایس دا مطلب اے کہ کوئی دِتّی ٹیورنگ مشین کسے دِتّے انپٹ
اُتے رُکے گی یا نہیں، ایہ حسابی طور اُتے ناقابلِ فیصلہ اے۔
ایس لئی اگلے قضیے نوں ``حساب رُکّن دے مسئلے دی ناقابلِ
فیصلہ حیثیت'' آکھیا جاندا اے۔ تھلّے اسی دو ثبوت دیاں گے۔
پہلا پچھلی بحث دی کڑی اگّے ودھاندا اے، دوجا ودھ سِدھا اے۔

\begin{thm}
\ollabel{thm:halting-problem}
من لو
\[
h(e, x)  =
\begin{cases}
1 & \text{جے\/ $\fn{Un}(e, x)$ تعریف شدہ ہووے} \\
0 & \text{ہور صورت وچ۔}
\end{cases}
\]
تاں $h$ قابلِ حساب نہیں۔
\end{thm}

\begin{proof}
فرض کرو $h$ قابلِ حساب اے۔ اسی وکھانا اے کہ ایس نال
اک ہمہ گیر قابلِ حساب فنکشن دی تعریف ہو سکدی اے۔ من لو
\[
\fn{Un'}(e,x) =
\begin{cases}
\fn{Un}(e,x) & \text{جے $h(e,x) = 1$} \\
0 & \text{ہور صورت وچ۔}
\end{cases}
\]
ہن $\fn{Un'}(e, x)$ اک کُل فنکشن اے، تے جے~$h$
قابلِ حساب اے تاں ایہ وی قابلِ حساب اے۔ مثال دے طور اُتے،
اسی بنیادی بازگشت نال $g$ دی تعریف کر سکدے آں:
\begin{align*}
g(0, e, x) & \simeq 0 \\
g(y+1, e, x) & \simeq \fn{Un}(e,x);
\end{align*}
فیر
\[
\fn{Un'}(e,x) \simeq g(h(e,x),e,x).
\]
$\fn{Un'}(e,x)$، $\fn{Un}(e,x)$ دے برابر اے
جتھے وی دوجا تعریف شدہ اے؛ ایس لئی $\fn{Un'}$
اوہناں جزوی قابلِ حساب فنکشناں لئی ہمہ گیر اے جیہڑے
اتفاق نال کُل نیں۔ پر ایہ \olref[nou]{thm:no-univ}
نال ٹکراندا اے۔
\end{proof}

\begin{proof}
فرض کرو $h(e,x)$ قابلِ حساب ہووے۔ $g$ دی تعریف ایویں کرو:
\[
g(x) =
\begin{cases}
  0                & \text{جے $h(x,x) = 0$} \\
  \fundefined & \text{ہور صورت وچ۔}
\end{cases}
\]
فنکشن $g$ جزوی قابلِ حساب اے۔ مثال دے طور اُتے، ایس نوں
$\umin{y}{h(x,x) = 0}$ دے طور اُتے تعریف کیتا جا سکدا اے۔
ایس لئی کسے~$e$ لئی $g(x) \simeq \fn{Un}(e,
x)$ ہر~$x$ اُتے۔ کیہ $g$، $e$ اُتے تعریف شدہ اے؟
جے اے، تاں $g$ دی تعریف مطابق $h(e,e) = 0$
(تعریف شدہ ہوئے بغیر~$h$ دی قدر~$0$ نہیں ہو سکدی)۔
$h$ دی تعریف مطابق ایس دا مطلب اے کہ $\fn{Un}(e, e)$
غیر تعریف شدہ اے۔ پر ساڈی فرض کیتی برابری
$g(x) \simeq \fn{Un}(e, x)$ ہر~$x$ لئی اے،
ایس لئی $g(e)$ وی غیر تعریف شدہ اے: تضاد۔
دوجے پاسے جے $g(e)$ غیر تعریف شدہ اے، تاں
$h(e,e) \neq 0$، تے ایس لئی $h(e,e) = 1$۔
ایس توں $\fn{Un}(e, e)$ تعریف شدہ ہوندا اے۔
پر $g(x) \simeq \fn{Un}(e, x)$ ہون کرکے $g(e)$
وی تعریف شدہ ہوندا۔ فیر تضاد۔
\end{proof}

\begin{tagblock}{TMs}
\begin{explain}
ایس دلیل نوں ٹیورنگ مشیناں دے لفظاں وچ وی بیان کر سکدے آں۔
فرض کرو کوئی ٹیورنگ مشین~$H$ ہووے جیہڑی ٹیورنگ مشین~$E$
دا بیان تے اک انپٹ~$x$ لے کے فیصلہ کرے کہ $E$ انپٹ~$x$
اُتے رُکدی اے یا نہیں۔ فیر اسی اک ہور ٹیورنگ مشین~$G$
بنا سکدے آں جیہڑی اکّو انپٹ~$x$ لے، $H$ نوں چلا کے
فیصلہ کرے کہ مشین $M_x$، جیہدا اشاریہ~$x$ اے، انپٹ~$x$
اُتے رُکدی اے یا نہیں، تے اوہدے اُلٹ کرے۔ ہور لفظاں وچ،
جے $H$ دسے کہ $M_x$ انپٹ~$x$ اُتے رُکدی اے، تاں $G$
لامتناہی چکر وچ پے جاندی اے؛ تے جے $H$ دسے کہ $M_x$
انپٹ~$x$ اُتے نہیں رُکدی، تاں $G$ بس رُک جاندی اے۔
کیہ $G$ اپنا اشاریہ انپٹ بنا کے رُکدی اے؟ اُتّے والی
دلیل دسدی اے کہ اوہ اُدوں تے صرف اُدوں رُکدی اے جدوں
اوہ نہیں رُکدی؛ ایہ تضاد اے۔ ایس لئی ساڈا ایہ فرض غلط
ہونا لازمی اے کہ ایداں دی کوئی ٹیورنگ مشین~$H$ موجود اے۔
\end{explain}
\end{tagblock}

\end{document}
